% !PRINTED 597
\weberpage{597}
Hiernach lässt sich die höchste Potenz des Primideals $\mathfrak{p}$ angeben, die in dem Partialgrundideal $\mathfrak{G}_R$ aufgeht. Nach \S.\,159 nämlich ist dies Ideal definirt durch das Product
\[ f'(\tau) = (\tau - \tau_1)(\tau - \tau_2)\ldots(\tau - \tau_{n-1}), \]
worin $\tau_1, \tau_2, \ldots, \tau_{n-1}$ die von $\tau$ verschiedenen unter den mit $\tau$ conjugirten Functionalen sind. Einer der Factoren dieses Productes, $\tau - \tau\mkern1mu|\mkern1mu\varphi$, ist dann und nur dann durch $\mathfrak{p}$ theilbar, wenn $\varphi$ zu $X$ gehört. Da nun die identische Substitution ausgeschlossen ist, so folgt, dass $\mathfrak{G}$ den Factor $\mathfrak{p}$ mindestens $g-1$ mal enthält.

Es ist aber $\tau - \tau\mkern1mu|\mkern1mu\chi$ dann und nur dann durch $\mathfrak{p}^2$ theilbar, wenn $\chi$ zu $X^{(1)}$ gehört, und folglich ist $\mathfrak{G}$ durch $\mathfrak{p}^{g-1+g_1-1}$ theilbar. Wir können so fortfahren und finden die genaue Potenz von $\mathfrak{p}$, die in $\mathfrak{G}$ enthalten ist. Sie hat den Exponenten
\[ (g-1) + (g_1-1) + (g_2-1) + \cdots \]
Jede der Zahlen $g, g_1, g_2, \ldots$ ist Theiler der vorangehenden, und $g_1, g_2, \ldots$ sind Potenzen von $p$, und die Reihe setzt sich so lange fort, bis eine der Zahlen $g = 1$ geworden ist. Die Zusammensetzung des Grundideals $\mathfrak{G}$ ist also vollkommen durch die Zerlegung der Gruppe $X$ bestimmt.

Die Zahlen $g, g_1, g_2, \ldots$ sind überdies für alle conjugirten Primideale $\mathfrak{p}_1, \mathfrak{p}_2, \ldots, \mathfrak{p}_e$ dieselben, und es folgt also
\[ \mathfrak{G} = \Pi(\mathfrak{p}_1\mathfrak{p}_2\ldots\mathfrak{p}_e)^{g-1+g_1-1+g_2-2\cdots}, \] % [sic] exponent pattern breaks from g_2-1 to g_2-2 as printed
wo sich das Productzeichen auf die Primfactoren aller Primzahlen $p$ des Körpers $\Omega$ bezieht. Es brauchen dabei aber selbstverständlich nur die in endlicher Anzahl vorhandenen Primzahlen $p$ berücksichtigt zu werden, für die $g > 1$ ist.

Die Partialnorm des Ideals $\mathfrak{G}$
\[ \mathfrak{D} = \mathfrak{N}_R(\mathfrak{G}) = \Pi\, \mathfrak{P}^{ef(g-1+g_1-1+g_2-1+\cdots)} \]
heisst die Partial-Discriminante des Körpers $\Omega$ in Bezug auf den Körper $R$. Sie ist ein im Körper $R$ gelegenes Ideal, die für den Fall, dass $R$ der Körper der rationalen Zahlen ist, von einem Einheitsfactor abgesehen, in die Grundzahl des Körpers $\Omega$ übergeht (\S.\,158).

\begin{center}
\S.\,163.

\gesperrt{Die Ideale in den Theilern des Körpers $\Omega$.}
\end{center}

Wenn die Gruppe $\Phi$ des Körpers $\Omega$ einen Theiler $\Phi'$ hat, so gehört zu $\Phi'$ ein Körper $\Omega'$, der ein Theiler von $\Omega$ ist und
%CONT>>
